% Generated draft: requires mathematical and source review.
\documentclass[11pt]{article}
\usepackage[margin=25mm]{geometry}
\usepackage{fontspec}
\setmainfont{Latin Modern Roman}
\usepackage{amsmath,amssymb,mathtools,unicode-math}
\setmathfont{Latin Modern Math}
\usepackage{enumitem,needspace,xurl,hyperref}
\setlength{\parindent}{0pt}
\setlength{\parskip}{5pt}
\setlength{\emergencystretch}{3em}
\newcommand{\N}{\mathbb{N}}
\newcommand{\Z}{\mathbb{Z}}
\newcommand{\Q}{\mathbb{Q}}
\newcommand{\R}{\mathbb{R}}
\newcommand{\C}{\mathbb{C}}
\newcommand{\F}{\mathbb{F}}
\newcommand{\A}{\mathbb{A}}
\newcommand{\PP}{\mathbb{P}}
\newcommand{\E}{\mathbb{E}}
\newcommand{\eps}{\varepsilon}
\newcommand{\dd}{\,\mathrm{d}}
\begin{document}
\section*{UnsolvedMath ID: 30003825}
\textbf{Source problem code:} catalogue entry
\subsection{Definitions and mathematical statement}
The ground field and its characteristic are not specified in the supplied record. The indexing of the basis of the positive Witt algebra is also not stated explicitly. Consequently, the following is a formulation schema requiring confirmation of these conventions, not a conjecture asserted uniformly over all fields.

Let \(k\) denote the intended ground field. Under the conventional positive-index interpretation, let
\[
W_+=\bigoplus_{i\geq 1}ke_i,
\qquad [e_i,e_j]=(j-i)e_{i+j}\quad(i,j\geq 1),
\]
where integer coefficients are interpreted in \(k\). Let \(A=U_k(W_+)\) be its universal enveloping algebra, namely the associative unital algebra generated by \(W_+\) subject to \(xy-yx=[x,y]\).

For a nonzero unital \(k\)-algebra \(B\), define its Gelfand--Kirillov dimension by
\[
\operatorname{GKdim}_k B
=\sup_{\substack{V\subseteq B\text{ a finite-dimensional }k\text{-subspace}\\1_B\in V}}
\limsup_{n\to\infty}\frac{\log\dim_k(V^n)}{\log n},
\]
where \(V^n\) is the linear span of products of \(n\) elements of \(V\).

Subject to the unresolved ground-field and indexing conventions, prove or disprove each of the following assertions:
\begin{enumerate}
\item For every nonzero two-sided ideal \(I\) of \(A\), the quotient \(A/I\) has finite Gelfand--Kirillov growth. Explicitly,
\[
\forall I\trianglelefteq A,\qquad
I\neq\{0\}\ \Longrightarrow\
\bigl(A/I=0\ \text{or}\ \operatorname{GKdim}_k(A/I)<\infty\bigr).
\]
The zero quotient is included as having bounded growth, so no properness hypothesis on \(I\) is added.
\item The two-sided ideals of \(A\) satisfy the ascending chain condition: for every sequence of two-sided ideals
\[
I_1\subseteq I_2\subseteq I_3\subseteq\cdots,
\]
there exists \(N\geq 1\) such that \(I_n=I_N\) for all \(n\geq N\).
\end{enumerate}
The main supplied statement quantifies over nonzero ideals. Its literature assessment instead says ``every nonzero quotient,'' which is a different condition: a nonzero quotient can have zero kernel, whereas a quotient by a nonzero ideal can be zero. The first assertion above follows the main statement; the intended literature wording requires review.
\subsection{Short English statement}
For the intended positive Witt algebra, ask whether every quotient by a nonzero two-sided ideal has finite Gelfand--Kirillov growth, and whether every increasing sequence of two-sided ideals eventually stabilizes. The ground field and basis indexing need confirmation. The supplied literature summary also uses a different quotient qualification.
\subsection{Sources}
\begin{itemize}
\item Local UnsolvedMath registry, \texttt{data/unsolvedmath/problems.json}, supplied record with permanent ID \(30003825\), title \emph{Ideals in the Positive Witt Algebra}, source problem code \texttt{catalogue entry}, in \emph{Oberwolfach Reports Open Problems}. No individual proposer is supplied; the proposed year is recorded as 2018.
\item Supplied report citation: \emph{Interactions between Algebraic Geometry and Noncommutative Algebra (2019)}, Oberwolfach Reports, pp.\ 1465--1515, DOI \texttt{10.4171/owr/2018/24}. Supplied source link and literature-source label ``original OWR report'': \url{https://doi.org/10.4171/owr/2018/24}. The record describes Oberwolfach Reports as documenting workshops and their open-problem discussions.
\item Supplied literature source, ``Sierra current research list'': \url{https://webhomes.maths.ed.ac.uk/~ssierra/preprints.html}.
\item Literature assessment dated exactly \texttt{2026-08-21}: the supplied curation labels the problem open, says the two conjectures were explicitly left unproved, and reports that later work located provides partial results rather than a resolution. This is supplied dated curation, not an independently verified assessment; its time-sensitive status and cited claims require independent verification.
\item The registry labels the statement \texttt{corrected\_verified} and reports a manual semantic, typographic, and LaTeX review. These are supplied provenance claims, not mathematical verification of this draft.
\end{itemize}
\end{document}
